Color each of $1,2,\dots,N$ with one of $r$ colors. Call a solution of $x+ y= z$ if $x$, $y$, and $z$ all receive the same color. ($x$ and $y$ need not be distinct.)

    
(a) Find a 2-coloring of $\{1,2,3,4\}$with no monochromatic solution, and prove that every 2-coloring of $\{1,2,3,4,5\}$
has one.


(b) Prove that for every r there is an $N(r)$ such that every $r$-coloring of $\{1,...,N\}$ with$ N \geq N(r)$ admits a
monochromatic solution. (Suggestion: color the edges of a complete graph on vertices 1,...,N+ 1 by giving the
edge $\{i,j\}$the color of $|i-j|$, and find a monochromatic triangle.)


(c) Here is the payoff, and it is not obvious that (b) has anything to do with it. Fix an exponent $k\geq 1$. Prove that
for every sufficiently large prime p, the congruence
$x^k + y^k \equiv z^k$ (mod $p$)
has a solution with $p\nmid  xyz$. (Suggestion: the nonzero residues mod $p$ split into $\gcd(k,p-1)$ classes according to
which coset of the $k$-th powers they lie in. How many classes can there be? Color by class.)


(d) Your argument in (b) gives a bound on $N(r)$. Compare it with the truth for $r = 2$, and find a 3-coloring of
$\{1,...,13\}$with no monochromatic solution. How far apart are the bound and the truth?
